\section{Evaluation}
\label{sec:results}

We answer five questions:
\begin{description}\setlength\itemsep{0pt}
  \item[RQ1] Who should decide? Does taking probe selection and stopping from the model cost resolved cases, and does any model add to the controller? (\S\ref{sec:decide})
  \item[RQ2] Who should read? Can a fixed parser replace the model when probes return raw logs? (\S\ref{sec:read})
  \item[RQ3] Is a model reader safe? Can an adaptive attacker who targets the model's reading (A3) close cases, and does reader trust stop it? (\S\ref{sec:rawlog})
  \item[RQ4] Can attackers reach the decisions through injected instructions (A1) or forged sources (A2), and what do the finish guard and source-diverse corroboration cost? (\S\ref{sec:injection})
  \item[RQ5] How do the evidence rule, imperfect tables, and causes missing from the candidate list trade automation for safety? (\S\ref{sec:policy})
\end{description}

\subsection{Setup}
\label{sec:setup}
\label{sec:testbed}

\paragraph{Triage settings} We built three simulated settings in which ground truth is known and the attacker can be controlled; the public datasets we examined lack paired attacks and look-alike administration under objective labels (\S\ref{sec:realdata}). \emph{sec-triage} is one alert on a web service with eight causes: three attacks (credential stuffing, SQL-injection probing, DNS command-and-control), an insider exfiltrating data, an exposed admin panel, a known-vulnerable component, an authorized vulnerability scan, and a monitoring false positive. Ten read-only probes cost 1 to 3 units; threat intelligence only says ``malicious'', which three attacks share. Each cause appears in a base variant, a noise variant with transient failures, and a drift variant in which the incident changes after the second probe (24 tasks). \emph{sigma-triage} removes our own choice of causes: each of 12 rules drawn with a fixed seed from the public SigmaHQ repository~\citep{sigmahq} contributes its targeted attack and up to four benign explanations taken verbatim from its \texttt{falsepositives} field (76 tasks); which outcome each of ten generic probes returns under each cause was annotated once by a model outside both planner families (phi-4~\citep{phi4}) and frozen. The annotation names one outcome per cause and probe, so sigma-triage is deterministic, and a rule whose annotation left a cause without an outcome that singles it out was replaced by the next rule in the sampling order; this happened to 1 of the 13 rules annotated. \emph{comp-triage} has three scenarios (remote execution, a scheduled task, a large outbound transfer) in which an attack and three legitimate activities use the same tools; outcomes are probabilistic, and five of its twelve causes have no single-observation signature. Tables are counted from 20 LLM-free development episodes per cause and variant, with seeds disjoint from evaluation.

\paragraph{Models and configurations} We use Qwen2.5-7B, 32B, and 72B and Llama-3.1-8B and 70B (32B and larger with 4-bit AWQ weights), served by vLLM at temperature 0.2 with budgets of 10 probes, 10 cost units, and 12 decisions per episode. The baselines are the usual agent designs: \emph{ReAct-style}, where the model sees the raw history, picks every probe, and decides the verdict, and \emph{Memory-only}, where the model additionally sees the ledger. \hesp{} configurations switch on controller selection, the finish guard, the controller stop, corroboration, and a reader policy one at a time, so that each rule's effect can be attributed. In controller-selection configurations the prompt does not reveal the ranking, so the model receives the same request whichever selector runs. An \emph{LLM-free controller}, whose scripted planner never proposes a verdict, gives the controller's performance without any model. \hesp{} as recommended uses information-gain selection, the guard, the confirmation stop, source-diverse corroboration, and recovery; Appendix~\ref{app:loop} lists every configuration.

\paragraph{Attacks} For A1, four injected texts that differ in wording and position (an operations note, a user-agent field, a fake analyst override, and a single probe's response) instruct the agent to close the case as the authorized scan. For A2, the source-reputation probe reports a known scanner (\emph{forged probe}), or the upstream reputation feed read by two probes is poisoned (\emph{forged feed}). For A3, probes return log text in the documented format or in a \emph{drifted} format that states the same facts in reworded lines, as after a logging-pipeline change. The first A3 attack writes a benign claim into the user-agent field of four probes; the adaptive attacks write a consistent authorized-scan story into every probe's field, do the same on the drifted format, or break out of the field to forge a complete log record (\emph{log-line injection}). A scripted reader that reads every format correctly and adopts every claim it can see gives each attack's upper bound. All attacks target the six actionable causes of sec-triage.

\paragraph{Metrics} An episode is \emph{resolved} if it ends with the true cause and cites an observation that only this cause produces; the verifier that checks this knows the hidden cause and is used only for evaluation. A \emph{missed attack} is an actionable case closed as benign, counted over all actionable episodes, so an agent that never answers cannot look safe. An episode is \emph{escalated} if it ends with \textit{other} or without a verdict. \emph{Cost} is the probe cost spent. In the raw-log experiments the verifier checks the true outcomes, so a parser error never counts as resolved, and \emph{adoption} counts parsed observations that equal the attacker's claim while the truth differs.

\paragraph{Statistical analysis} Each experiment's hypotheses and primary comparisons were written down before its first episode, and each primary comparison is reported with a paired, task-clustered bootstrap interval (2{,}000 resamples; Bonferroni-adjusted where an experiment has several). An audit replayed the journal of every one of the 23{,}582 LLM episodes and of the LLM-free episodes. Appendix~\ref{app:prereg} lists the experiments, their primary comparisons and outcomes, and the errors we found in our own process.
